\begin{tikzpicture}[font=\figurefont\small,
 box/.style={draw=BookRule,fill=BookPaper,minimum width=17mm,
 minimum height=8mm,align=center},
 flow/.style={-{Latex},draw=BookInk,line width=.8pt},
 grad/.style={-{Latex},draw=BookAmber,dashed,line width=.8pt}]
 \foreach \y/\tag in {0/VDN,-3.4/QMIX} {
   \node[anchor=west,text=BookTeal] at (-.8,\y+.8) {\tag};
   \node[box,fill=FigureInput!12] (h1\tag) at (0,\y) {$h_1$};
   \node[box,fill=FigureInput!12] (h2\tag) at (0,\y-1.2) {$h_2$};
   \node[box,fill=FigureModel!12] (q1\tag) at (2.5,\y) {$q_1(h_1,\cdot)$};
   \node[box,fill=FigureModel!12] (q2\tag) at (2.5,\y-1.2) {$q_2(h_2,\cdot)$};
   \node[box,fill=FigureOutput!12] (out\tag) at (8,\y-.6) {$Q_{\text{tot}}$};
   \draw[flow] (h1\tag)--(q1\tag);
   \draw[flow] (h2\tag)--(q2\tag);
 }
 \node[box] (sum) at (5.4,-.6) {$q_1+q_2$};
 \draw[flow] (q1VDN)--(sum);
 \draw[flow] (q2VDN)--(sum);
 \draw[flow] (sum)--(outVDN);
 \node[box] (mix) at (5.4,-4) {$f_s(q_1,q_2)$\\逐坐标非减};
 \node[box,fill=FigureInput!12] (state) at (5.4,-2.4) {训练状态$s$};
 \draw[flow] (state)--(mix);
 \draw[flow] (q1QMIX)--(mix);
 \draw[flow] (q2QMIX)--(mix);
 \draw[flow] (mix)--(outQMIX);
 \draw[grad] (outVDN.south)--(8,-1.9)--(2.5,-1.9)--(q2VDN.south);
 \draw[grad] (outQMIX.south)--(8,-5.3)--(2.5,-5.3)--(q2QMIX.south);
 \node[text=BookMuted,anchor=west] at (-.8,-6)
   {局部执行：各自取$\arg\max q_i$；虚线代表对各效用的训练反传。};
\end{tikzpicture}
